\chapter{从数学模型到可维护实现}
\section{数据层保持科学约定}
程序以四个核心对象组织计算。\code{Cell} 保存晶胞和度量，\code{Atom} 保存物理占有率、分数坐标和位移张量，\code{Structure} 保存原子与对称操作，\code{Reflections} 保存观测、标准差、批次和留出标志。对象建立时检查维度、有限数值、元素、唯一原子标签和物理范围。

内部约定稳定后，CIF、SHELX 与 JSON 之间的差异集中在输入输出边界。JSON 保存完整浮点精度及明确字段。传统 HKLF4 的固定列宽可能限制输出精度，因而精确往返使用 JSON；HKLF4 导出不会暗中把无法表达的留出标志丢失。

\section{模块对应的数学职责}
\begin{longtable}{p{43mm}p{93mm}}
\toprule 模块 & 数学或执行职责\\\midrule\endfirsthead
\toprule 模块 & 数学或执行职责\\\midrule\endhead
\code{model.py} & 晶胞矩阵、度量、周期最近像、物理数据对象\\
\code{symmetry.py} & 空间群操作、特殊位置、消光与候选诊断\\
\code{scattering.py} & 原子散射因子、对称展开、热位移、非相干畴强度\\
\code{phasing.py} & 带随机种子的 $P1$ 电荷翻转与稀疏原子投影\\
\code{refinement.py} & 参数化、信赖域估计、冻结权重、秩与协方差\\
\code{restraints.py} & 几何和位移限制残差\\
\code{absolute.py} & 显式异常散射因子与反射对商量分析\\
\code{maps.py} & 统一 Fourier 约定的密度、差值与 Patterson 图\\
\code{visualization.py} & 概率椭球、等值面、倒易空间、矢量图与离线交互图\\
\code{validation.py} & 反射轨道留出与观测诊断\\
\code{feedback.py} & 有信息上限的结构先验轮廓分解\\
\code{io.py}, \code{shelx.py} & 结构与反射格式转换、原始指令的无损保存\\
\code{compat.py} & 单独安装的官方程序执行、附属文件与运行记录\\
\code{agent.py}, \code{cli.py} & 结构化动作、输入快照、运行目录与机器结果\\
\code{mcp_server.py} & 官方 MCP SDK 适配层\\
\bottomrule
\end{longtable}
模块之间通过显式对象和返回值通信。相位恢复不调用语言模型，精修不修改原始反射，绘图不决定删点，兼容层不偷偷代替原生算法。该边界使单个数学部件可以独立测试与替换。

\section{复杂度与实际优化方向}
设独立反射数为 $H$，展开后的原子位置数为 $N$。直接结构因子求和的主要计算量为 $O(HN)$。对于 $P$ 个独立参数，三点差分构造完整 Jacobian 约需 $2P$ 次扰动计算，故一次完整导数评估的量级为 $O(PHN)$。密集法矩阵的分解约为 $O(P^3)$。

这说明大型 MOF 的性能瓶颈不能单靠外层 agent 并行解决。直接有效的优化是缓存不随参数变化的晶胞与散射数组、实现解析或自动微分 Jacobian、利用约束稀疏性，并在独立相位试验或候选模型之间并行。GPU 可以加速数组运算，但主机与设备的数据复制、动态特殊位置处理和不必要的全矩阵构造仍会限制性能。

配套实现采用 NumPy 向量化、SciPy 求解和 Gemmi 表格及文件基础设施。它没有声称实现 GPU、稀疏 CGLS 或全解析 Jacobian，也没有用不同数据或不同收敛准则给出与官方程序的速度排名。性能评价应在相同硬件、相同数据、相同模型和相同接受标准下进行。

对体素总数 $M=N_1N_2N_3$，一次三维 FFT 的复杂度约为 $O(M\log M)$，多起点相位恢复的计算量与试验数和迭代数近似成正比。相位试验之间相互独立，适于并行；同一试验中的迭代则按前一步结果顺序推进。

\section{验证层次}
验证应逐层建立独立性。闭式量检验包括晶胞度量、单参数尺度和概率椭球；物理不变量包括原点平移强度不变、Friedel 对称及其异常散射破缺、对称等价反射、特殊位置多重性和张量变换；外部数值对照使用 Gemmi 的独立结构因子计算；工作流验证检查输入不变、求解后精修、真实数据回归与工具传输。

合成数据由被测正向模型生成时，反演恢复能够检验工作流和优化，但不能独自证明正向模型正确。因此，需要外部实现、解析结果与真实观测三种证据互相补充。附录明确列出每项验证的来源与所能支持的结论。

\section{完整功能迁移的验收原则}
SHELXT 与 SHELXL 的功能由大量文件约定、默认值和数值分支共同组成。重新实现某个公式，不等于复现某条完整指令。一个可靠迁移单元应同时给出：输入语法与默认值；参数的数学定义；原生执行结果；异常和边界行为；与公开归档或官方程序的对照证据。

因此，本项目为官方指令建立逐项状态，而不是一个总体“兼容”布尔值。当前原生范围由源码与测试确定；原始指令文件保持不变并可交给官方程序。以后某个功能完成独立实现和对照测试后，才能将其从外部执行迁入原生能力表。

结构精修的成熟还需要跨化学体系的数据集，包括轻原子、重原子、手性、低完整度、强吸收、伪对称、无序、不同孪晶类型和大晶胞。对这些任务的覆盖，应以可复现样本和失败模式积累，而不是把一次成功求解外推为普遍等价性。
